\documentclass[11pt]{article}
\usepackage[T1]{fontenc}
\usepackage{lmodern}
\usepackage[margin=1.25in]{geometry}
\usepackage{setspace}
\setstretch{1.12}
\usepackage[colorlinks=true,linkcolor=black,citecolor=black,urlcolor=black]{hyperref}

\title{Magnificent Monotony and the Non-Player Character:\\Swotting, Grinding, and the Supervision of Purpose}
\author{Flyxion\\Independent Researcher}
\date{November 2026}

\begin{document}
\maketitle

\begin{abstract}
\noindent A passage in the posthumous notes collected as \emph{The Will to Power} names the classical philologist as the educator par excellence, on the ground that his work offers the best pattern of magnificent monotony in action and that under his banner the young learn to swot. The passage describes in the vocabulary of nineteenth-century schooling what contemporary players call grinding: the repeated performance of a low-reward task for the sake of a later capacity. This essay develops the resemblance into an argument. It holds that Nietzsche's observation that nothing is more rare than a personal act is best read as a claim about the relation of an agent to its purposes and not about behavioral richness, that the figure of the non-player character captures this relation more exactly than the figure of the unconscious automaton, and that the relation cannot be read from the routine itself. It offers a modular account of why such regimens persist: repetition reduces the deliberative cost of routine operations, which can free higher-order judgment or leave it subordinated to the routine, so that automaticity and autonomy are not opposites. It relates the criterion to Nietzsche's rejection of requital, arguing that a system of equivalence presupposes interchangeable role-occupants and that responsiveness permits contributions a role's specification cannot fully capture, without treating uniqueness as a condition of autonomy. The essay separates three achievements that are easily run together, originating a purpose, endorsing it, and remaining able to revise it, and argues that the third is the one that matters. It proposes responsiveness as the working criterion, defined operationally as sensitivity to relevant changes in evidence, consequences, and available alternatives, where the response may be persistence, modification, or abandonment. It then separates the question of authorship from the question of what an activity does in the world, since an independently chosen activity may destroy what an assigned one preserves.
\end{abstract}

\newpage
\section{Two Vocabularies for One Structure}

Two vocabularies describe the same structure, and each hides a feature that the other makes plain. The first belongs to a philologist who became a philosopher and who, in notes he never arranged for publication, set down a remark about the profession that had formed him. The remark appears in the Ludovici translation as follows.

\begin{quote}
It is on this account that the philologist has, hitherto, been the educator per se: because his activity, in itself, affords the best pattern of magnificent monotony in action; under his banner youths learn to ``swat'': first prerequisite for the thorough fulfilment of mechanical duties in the future (as State officials, husbands, slaves of the desk, newspaper readers, and soldiers). \cite[Aph.\ 888]{wtp2}
\end{quote}

The second vocabulary belongs to the players of online role-playing games, who use the verb to grind for the repeated performance of a task whose immediate reward is small, undertaken because repetition is the price of a later capacity, rank, or possession. The two descriptions concern different centuries and different institutions, yet they pick out the same arrangement: an activity whose individual instances are of little interest, whose value is deferred, and whose main effect on the performer is to produce a disposition to continue. The word Ludovici sets in quotation marks as swat is a variant of the British swot, which means to study laboriously and by repetition, so that the older spelling, set off in quotation marks, already names the grind.

The resemblance would be a curiosity if it stopped at vocabulary. It does not, because the passage invites a question that the game vocabulary makes precise: what separates repetition that expresses an independent character from repetition that merely discharges a role? The sections below argue that the answer concerns the relation of the performer to the purpose the repetition serves, that this relation has more than one component, and that it must be kept distinct from a second question, namely what the repetition accomplishes in the world. A reading of the contemporary grind as a form of labor is available in the literature on game studies \cite{yee2006}, and the convergence suggests that something general is being described.

\section{The Philologist as Educator}

The claim that the philologist's activity is valuable \emph{in itself} as a pattern of monotony is a claim about form and not about content. The passage does not say that the study of Greek grammar improves the mind through its subject matter, nor that the texts transmit a worthy culture. It says that the activity, apart from what it is about, offers a pattern, and that the pattern is what the pupil imitates: a posture of continuing without being moved by the object, of applying effort where no personal interest sustains it. The aphorism says as much in its own words: the task of the higher schools is to learn something one does not care a fig about and to find one's duty in that objective activity \cite[Aph.\ 888]{wtp2}.

The biographical context supports a skeptical reading of the institution. Nietzsche was appointed to a chair of classical philology at Basel in 1869 and left it in 1879, so that the remark comes from an insider assessing the function of his own training. The earlier notes published as \emph{We Philologists} are also posthumous and unfinished, and they are cited here only to show that doubt about the philologist's pedagogy was long-standing; the translator's synopsis lists among their themes the philologist's inability to train his pupils \cite{wp}. It is natural to read magnificent monotony in this light, as a sardonic description of a large, durable, and socially valuable regime whose educational yield is endurance and not understanding.

The surrounding aphorism settles the function. It opens by announcing an economic justification of virtue whose object is to make man approximate an infallible machine, equipped with machine-like virtues, and it names the tedium and monotony of mechanical activity as the first stumbling-block. Learning to endure that tedium is described as the task of all higher schools, and the philologist is the educator per se because his work is the best pattern of such monotony and the first prerequisite for mechanical duties as official, husband, clerk, newspaper reader, and soldier. The tone is less settled. The aphorism closes on a mechanical form of existence regarded as the highest form and worshipping itself, which suggests critical distance, but it opens in the first person with a project announced and not condemned. The adjective may be mocking, but it may also be descriptive, and the passage does not decide. The ambiguity matters because Nietzsche elsewhere insists, in a work he published, that anything of lasting worth on earth requires obedience over a long time in a single direction \cite[Sec.\ 188]{bge}, and he treats the constraining discipline of language, custom, and craft as the condition under which freedom of spirit becomes possible. The philologist may therefore be at once an object of criticism, as the producer of docile subordinates, and an exemplar of a capacity that Nietzsche regards as indispensable. This essay does not resolve the ambiguity. It takes it as the problem: the same discipline is a precondition of high achievement and a preferred instrument of systems that wish to produce reliable subordinates, and the evaluation depends on what the discipline serves and on the standing of the one who undergoes it.

A caution about the source is also required. \emph{The Will to Power} is not a book that Nietzsche wrote. It is a posthumous compilation of notebook material, arranged by editors under a plan that Nietzsche had abandoned, and its numbering and grouping are editorial constructions \cite{montinari2003}. Two aphorisms that stand in the same volume may come from different notebooks, years, and purposes, and their proximity in the compilation is not evidence that they belong to one sustained argument. The aphorisms on which this essay leans most heavily, 886 to 888, are adjacent in Book IV, which makes a joint reading natural, but the adjacency is the editors' and has not been checked against the notebooks. The remarks used here are therefore treated as independent observations by the same author, each of which can be assessed on its own, and not as steps in a single design. Where the argument below joins them, the joining is the essay's own.

\section{Mechanical Activity as Remedy}

A second passage, in the first volume of the same compilation, adds a function to monotony. Among measures grouped under the hygiene of the weak, the text lists as relieving measures absolute obedience, mechanical activity, and total isolation from men and things that might exact immediate decisions and actions \cite[Aph.\ 45]{wtp1}. The context is therapeutic. The person addressed reacts too readily to every stimulus and cannot refrain from responding, so that each demand to decide is a depletion. For such a person mechanical activity is a remedy because it removes the demand: the rules are fixed in advance, each step follows from the last, and the performer can proceed without consulting the competing impulses that would otherwise exhaust him.

Set beside the philologist passage, which is a separate observation, this one shows that a single structure can serve different functions depending on who stands in front of it. Applied to the young and undeveloped, repetitive tasks may form a disposition to endure. Applied to the exhausted, they supply relief from deliberation. Applied to the competent, they can serve as the practice through which a skill is acquired. The structure of the activity, a sequence of acts whose rules are given and whose outcomes are predictable, is the same in all three cases, and this is what makes it impossible to read the function from the form. The functions differ, moreover, in what they do to responsiveness itself. Training that builds the disposition to continue shapes the performer so that continuation becomes easier and less dependent on interest, and in doing so may reduce the occasions on which the question of ground arises. Relief suppresses the demand to reconsider for as long as it lasts and leaves the performer's capacity as it was. Practice by the competent leaves the question of ground open and serves it. The first alters the disposition, the second postpones the exercise, and the third employs the exercise. A person who spends a winter laying tile, a student who conjugates irregular verbs, and a player who farms the same dungeon for the hundredth time may each be doing something of very different significance, and an observer who attends only to what the hands are doing will not be able to say which.

\section{The Anatomy of the Grind}

The player's notion of the grind has a more precise structure than its casual use suggests. A grind is a loop in which an action is repeated, each repetition yields a small increment of a measured quantity, and the accumulated quantity unlocks something the player wants. Three features follow. The increments are individually unrewarding, which distinguishes grinding from play in the sense in which a game is the voluntary attempt to overcome unnecessary obstacles \cite{suits1978}: the grind retains the voluntary form of that attempt while draining out the interest in overcoming. The measure is external to the activity, in that the quantity gained is defined by the designer and not by the player. And the unlocked reward is specified in advance by the same designer, so that the grind both sets the price and fixes the good that the price purchases.

The student who swots does so under an educational system that fixes the curriculum, the examination, and the credential. The player who grinds does so under a game that fixes the loop, the meter, and the unlock. In both cases the performer is free in the sense that the activity may be abandoned, and in both cases the freedom is exercised within a space whose ends were laid down elsewhere. A person who repeats a scale for a year in order to play a piece of their own choosing, or who sands a hundred joints in order to build a house of their own design, is also grinding in the formal sense, but the meter and the reward belong to the performer. The activity has the same texture and a different owner.

The distinction between these cases is not between enjoyed and unenjoyed repetition. The designed grind can be pleasant, and its design often aims at making the repetition pleasant enough to continue, while the self-owned grind can be miserable throughout. Nor is it between the useful and the useless, since the designed grind frequently produces real capacities. It concerns the performer's relation to the end that gives the repetition its point, and that relation requires more careful analysis than the word authorship usually receives.

\section{Monotony as Modularity}

A further account of why such regimens persist treats repetition as a social discovery of something like modularity. In software, a module is valuable to the degree that it hides its internal decisions behind a narrow interface \cite{parnas1972}, and a function whose result depends only on its arguments can be understood, tested, and reused without reference to the program around it \cite{hughes1989}. Repetition does something similar in a human performer. A procedure practiced until it runs without deliberation is cheap to execute, hard to disturb, and easy to combine with others, and the saving in deliberative effort is what makes mastery of a narrow task attainable. The defensible claim is narrower: repetition can reduce the deliberative burden of certain operations and thereby free attention for other activities. It does not follow that tedium is the mark of this process. A skill can become automatic without being tedious, and a repetitive activity can remain tedious precisely because it has not become automatic. Nor does it follow that expertise excludes context. The experienced philologist, musician, or carpenter performs familiar operations effortlessly while growing more sensitive to distinctions that a novice cannot perceive.

The software analogy is instructive only if two things are kept apart. Encapsulation makes internal procedures reliable and independently maintainable by hiding implementation decisions. A well-encapsulated module may nonetheless remain fully responsive to changing inputs, environmental conditions, and explicit error states. Insulation is different: it prevents relevant information from reaching the procedure at all. Encapsulation can support agency, because it lets a performer rely on what has been mastered while attention goes elsewhere. Insulation can undermine it, because a procedure that cannot be reached by relevant information will continue after its point has gone. Information hiding is therefore not context independence, and the question that matters about any trained routine is which of the two it has.

Instruction in language offers a familiar contrast. A drill that isolates a grammatical form and repeats it until it is produced without thought builds procedural reliability, whereas reading or conversation in which the same form appears among many others builds the ability to notice when a form fits and when it does not. The two are not rivals. The first lowers the cost of the form and the second keeps the form answerable to what is being said, and a learner who has only the first can produce correct sentences that miss the point. The contrast illustrates the general difference between procedural reliability and contextual responsiveness without requiring any claim about what educators historically believed.

The consequence for the relation of an agent to its purposes is that automaticity and autonomy are not opposites. Automaticity can make autonomy possible by reducing the cognitive cost of routine operations, so that judgment is spent where it is needed. The danger arises when an agent's higher-order judgment becomes subordinated to the continuation of the operations, and not when the operations are automatic. Whether attention freed by repetition is redirected toward higher-order judgment or simply captured by another routine is a separate question that the regimen alone does not settle. This also explains why decades of repetitive study can produce intellectual independence rather than undermine it, and it does so without supposing that self-directed practice differs in its mechanics from externally assigned practice. The difference lies in whether a supervisory capacity retains authority over how and whether the routine is applied, a capacity the next sections try to specify.

\section{Three Achievements}

Talk of who authors a purpose runs together at least three different things. The first is originating a purpose: being the first source of it, so that it arises from the agent and from no one else. The second is endorsing a purpose: taking it up after reflection as one's own, whatever its origin, so that the agent would affirm it on examination. The third is remaining able to revise a purpose: retaining the capacity to reconsider the relation between the purpose and the activity when the circumstances supporting that relation change. These are different achievements and they come apart. A purpose can be originated by an agent who never examines it, endorsed by an agent who did not originate it, and revisable by an agent who has done neither with much care.

Origination is the least important of the three, and it is also nearly unattainable. Human purposes are ordinarily inherited, learned, negotiated, and shared. A person who wants to cure disease, build houses, or play the cello wants something supplied by a culture that existed before them, and the social origin of the want does not make it inauthentic. A requirement that every worthwhile purpose originate entirely within the individual would convict everyone, including those Nietzsche regarded as exemplary, who plainly took over inherited problems and traditions and made them their own. What matters is the capacity to evaluate and appropriate what one has been given.

The analogue in the philosophy of action is the distinction between persons and what Frankfurt called wantons \cite{frankfurt1971}. A wanton has desires and acts on them, sometimes with great skill, but does not care which desire prevails. A person has second-order attitudes toward their desires and can identify with some and reject others. This is a form of endorsement. Kant's account of immaturity as the inability to use one's understanding without the guidance of another, self-incurred when its cause is a lack of resolution and not of understanding \cite{kant1784}, points to the third achievement: the willingness and ability to think the matter through when guidance is available and the question is open. In each case the line falls between competence and a certain relation to one's ends, and in each case a person can be highly competent while lacking that relation.

On this reading, which is a philosophical reconstruction and not a settled historical interpretation (the status of the compilation is discussed in \cite{montinari2003}), Nietzsche's observation that nothing is more rare than a personal act \cite[Aph.\ 886(a)]{wtp2} is a claim about the second and third achievements. The same aphorism adds that class, rank, race, environment, and accident are far more likely to be expressed in a deed than the personality of the doer, so that most acts are legible as outputs of a position and not of a person. It adds that the majority are not even one personality, that from the standpoint of the species a personality would be an unnecessary expense and a luxury, and that in such circumstances everybody is a channel or a transmitting vessel \cite[Aph.\ 886(b)]{wtp2}. In the ordinary modern sense everyone has a personality, a configuration of traits and habits. In the stronger sense the observation requires, a personality is a human being whose drives have been organized under a coherent direction that the person has examined and taken up. To say that most people lack one is not to say that they lack intelligence, feeling, or skill. It is to say that the governing purposes of their conduct were supplied by convention, role, and institution and have not been examined and reorganized. The claim is severe and has an elitist tone, but its content is structural.

\section{The Non-Player Character}

The figure of the non-player character captures this structure with some economy. In game design the term names a character whose conduct is determined by the designer and not by a player. What distinguishes such a character is not simplicity of behavior. The inhabitants of a well-made world keep schedules, respond to the weather, remember what the player has done, and may pursue elaborate routines that no player observes. Behavioral richness does not separate the non-player character from the protagonist. What separates them is that the objectives of the former are specified from outside and are not available to it for revision. The character does not choose to keep the shop; keeping the shop is what the character is for.

If personality is a certain relation of the agent to its purposes, then its absence is not poverty of behavior but fixity of purpose, and an agent can have an enormous repertoire of competent behaviors while remaining, in this sense, a non-player character in a world whose author is society, the school, the employer, or the conventions of a class. The philologist passage then describes a mechanism of production. Under a regime of magnificent monotony the young can acquire the capacity to continue acting in the absence of personal interest, and this is the capacity a system needs from agents whose purposes it supplies. An agent who can act only when interested requires persuasion at each step, whereas one who can act regardless of interest can be directed.

The first book of the compilation supplies a figure of the same type in the virtuous man, who is said to be of a lower species because he has no personality and acquires his value by conforming with a scheme fixed once and for all \cite[Aph.\ 319]{wtp1}. The virtuous man is a non-player character in the exact sense: his value is conferred by the scheme he fits and not by anything he authored, and another occupant of the position would be worth the same. The remark about the rarity of the personal act generalizes the point, since it implies that most conduct, however skilled or worthy, is better explained by the role it discharges than by the person who discharges it.

The analogy also exposes an inversion. Within the game the player is the protagonist and the non-player characters are the dependents, and the player chooses to grind. But the grind is a loop designed by someone else, and the unlock that rewards it is chosen by someone else, so that relative to the designer the diligent player occupies the position that the non-player character occupies relative to the player. The same agent can be protagonist at one level and executor at another, and nothing in the experience of choosing announces which level one is on. A player who has ground for a month to obtain an item that the game made desirable has exercised choice at every step and has revised nothing. The cases below make this structure concrete.

It follows that the figure does not support the cheap use to which it is often put. In contemporary discourse the term has become an insult directed at persons whose opinions are thought to be received and not reasoned, with the suggestion that some human beings are less than fully human. Nothing in the distinction requires this. It concerns the organization of purposes, which is a matter of degree, of circumstance, and of effort, and which varies within one life across domains. The same person may revise their work endlessly and merely execute their civic habits, or the reverse. The proper use of the figure is as a diagnostic of a structure of dependence and not as a classification of people.

\section{Responsiveness to Reasons}

If form does not reveal function and behavior does not reveal the agent's relation to its purposes, what does? The proposal of this essay is that the relevant difference is modal and not occurrent. It lies in how the agent would respond when the circumstances that support the relation between an activity and its purpose change, and not in what the activity is doing now.

A tempting formulation is that the independent agent is the one who would abandon a routine when its purpose disappears. This is too simple, and it is wrong in both directions. A person who persists after an apparent purpose has gone is not necessarily dependent. They may be preserving a possibility that others have not recognized, pursuing an objective with a longer horizon than the apparent one, or refusing an institutional judgment that the project is finished. Conversely, a person who abandons a project the moment its external rewards vanish may be showing dependence on those rewards and not independence from them. Readiness to abandon is not the criterion, since sometimes the right response to changed circumstances is to continue.

The stronger criterion is responsiveness to reasons. What matters is whether the agent can reconsider the relation between the activity and its purpose when the circumstances supporting that relation change, and whether the outcome of the reconsideration tracks the considerations that bear on it. Sometimes reconsideration warrants abandonment, sometimes persistence, and sometimes a reorganization of the purpose itself. Two people performing the same grind can differ entirely in this respect and show no difference until the ground shifts. The one for whom the repetition was a means held under standing judgment will notice the shift, weigh it, and respond for reasons. The one for whom the repetition was the thing that relieved them of asking what to do next will continue, or will seek the next externally supplied loop, without the question ever being raised.

Responsiveness is defined operationally as sensitivity to relevant changes in evidence, in consequences, and in the alternatives available, shown in how conduct changes or fails to change when those things change. It is not defined by the ability to give reasons. An agent can produce elaborate justifications for continuing an activity while remaining wholly dependent on its externally imposed objectives, and such rationalization is distinguished from responsiveness by the fact that it would be offered whatever the evidence, consequences, or alternatives turned out to be. Conversely, a practically skilled agent may have excellent grounds for persisting and lack the verbal resources to state them, yet show them in what they attend to, what they test, and what would make them stop. The criterion counts practical competence as a form of responsiveness, shown in what the agent attends to, tests, and would stop for, and the response it looks for may equally be persistence, modification, or abandonment.

The criterion retains a residual circularity that should be owned. Which changes count as relevant is the very thing responsiveness is supposed to settle, and an agent who ignores a change as irrelevant is either judging correctly or insulating itself against a relevant input, which the definition alone does not distinguish. Responsiveness is accordingly a virtue-term and not a decision procedure. It does not generate judgments from scratch; it organizes retrospective and comparative ones, in which the later course of events, or the conduct of others in like circumstances, shows what was relevant.

In the terms of the modular account, responsiveness is the maintained openness of a routine's interfaces to relevant inputs, and insulation is its failure mode. The supervisory capacity is not best understood as a further module placed above the routine, since such a module would face the same question about its own inputs. It is a standing disposition of the performer to keep the routine answerable to what bears on it, and it is exercised in the routine's operation and not beside it.

A qualification drawn from the source prevents a misreading of the criterion as a demand for general openness. In the same aphorism a personality is described as a relatively isolated phenomenon, which requires timely isolation, something like a walled enclosure, a capacity for shutting out the world, and a much lower degree of sensitiveness than the average person, who is too easily infected with the views of others \cite[Aph.\ 886(c)]{wtp2}. This does not contradict the criterion. Responsiveness as defined here is sensitivity to what bears on the purposes at issue, evidence, consequences, and alternatives, and not susceptibility to the opinions of those nearby. A well-bounded agent is selectively permeable: closed to contagion and open to relevant considerations. In the terms of the modular account, susceptibility is a failure of the boundary in one direction and insulation a failure in the other. The word supervision also appears in the source in the opposite sense. In the next aphorism the strongest natures are said to be tightly bound, watched, and supervised by the instinct of the herd, under a regime of duties in which one can no longer call one's soul one's own \cite[Aph.\ 887]{wtp2}. The supervision of this essay is the internal counterpart of that external one, and the difference between them is the difference between a routine watched from outside and a routine kept answerable from within.

The criterion explains several features of the earlier discussion. It explains why the hygiene of the weak is a remedy and not a cure, since relief from deliberation is by construction relief from the exercise that would build responsiveness. It also shows why the philologist's regimen is attractive to a system that supplies purposes: a pupil trained only to continue regardless of interest may become insensitive to the failure of ground, and insensitivity of that kind makes an agent reliable to whoever supplies the purposes. This is an institutional outcome and not a necessary consequence of repetition. Sustained philological work can equally cultivate the critical capacities needed to question inherited interpretations, as Nietzsche's own passage from the discipline to its criticism suggests, and the passage establishes that the discipline trains endurance without establishing which of the two outcomes follows. And it explains the praise of long obedience. Obedience over a long time in a single direction is the means by which a purpose that has been examined and taken up becomes effective, because the world resists and effect requires persistence. What distinguishes it from the swotting of the philologist's pupil is that the persistence is held by a standing judgment that remains open to the reasons that would end or alter it, and not by habit alone.

The criterion is also a standard by which institutions can be assessed without admiring or despising discipline as such. A school, a workplace, or a game can be examined for whether it cultivates responsiveness to reasons about the purposes it asks people to serve, by teaching the grounds of the regimen, exposing them to inspection, and rewarding the reconsideration of tasks whose point has changed. The contrast between drill and integrated use supplies the model: the drill lowers the cost of the form, the integrated use keeps the form answerable to what is being done with it, and an institution that aims at responsiveness provides both and provides them in sequence. An institution that cultivates only endurance, and treats the question of ground as an impertinence, tends to produce the disposition the philologist passage describes, whatever it professes.

The comparison with the game is meant to isolate a structure of dependence, not to equate the stakes, and the institutional claim is qualified by the conditions under which the structure operates. The analogy between schooling and the game is imperfect in a respect that matters: a student cannot generally replace the school with another as a player can replace one game with another, and the costs of exit differ in kind. Where exit is costly, as in compulsory schooling or in work on which a livelihood depends, the structure operates with greater force and the case for cultivating responsiveness is correspondingly stronger.

\section{Fungibility and Requital}

The figure of the non-player character has a counterpart in the logic of exchange, and the criterion of responsiveness bears on it. Two aphorisms in Book IV approach the same structure from the side of requital. The first is a note on the golden rule, which it reads as the formula of a type of man in whom the instinct of the herd speaks, since the maxim presupposes that an action is invariably requited and that actions are equivalent. The note denies the equivalence: among real individuals equal actions do not exist, so that there can be no such thing as requital, and nobody can pay the speaker back for what he does, since the most that can be done is to make him the victim of another action \cite[Aph.\ 925]{wtp2}. The second, directed against Mill, objects to founding human traffic upon mutual services, so that every action appears as a cash payment for something done to us. It calls the hypothesis of equivalence ignoble, since it cancels the part of an action that has no equivalent and cannot be remunerated \cite[Aph.\ 926]{wtp2}. The same aphorism allows that equivalence holds in the very select circles where one actually has one's equal, and the concession matters: the claim is not that all exchange is illegitimate, but that exchange presupposes interchangeability and that interchangeability is not a fact about persons in general.

The point connects directly with the modular picture. A module is defined by its interface, and any instance that meets the specification can replace any other, so that a role-bound executor is fungible by design. The non-player character is the clearest case: another instance running the same routine would be the same shopkeeper, and nothing would be owed or lost in the substitution. The virtuous man of the first book is in the same position, since he has no independent value, may be compared, and has his equals \cite[Aph.\ 319]{wtp1}. The regimens discussed above tend to produce this fungibility in the persons they train. A credential, a rank, or a quantity of experience points is an equivalence ledger that converts so much effort into so much standing, and it works only if the efforts are comparable and the performers interchangeable. The grind is accordingly a small requital machine, and the meter that governs it is the device by which the performer's contribution is made equal to anyone else's.

Non-fungibility should not be confused with contingency, and it does not follow from responsiveness alone. Every act is unrepeatable in the trivial sense that no event occurs twice, and an executor of a role performs unrepeatable events constantly without any of them being irreplaceable. Nor is uniqueness of outcome the point. Two agents who have each examined their purposes and responded properly to the evidence may reach the same decision, and an institution that cultivates good judgment may produce considerable convergence among its members. Conversely, a person whose conduct is driven by an idiosyncratic habit may be practically irreplaceable without exercising any supervision over that habit. Uniqueness is therefore neither necessary nor sufficient for the relation the essay describes. The stronger proposition is that responsiveness permits a form of individual contribution that a role's existing requirements cannot completely specify. An agent who supervises a routine can revise the terms of performance and not merely satisfy them, and what a ledger of equivalents cannot register is this capacity to alter what counts as the performance. That preserves the connection with Nietzsche's rejection of requital, since a contribution that revises the terms of exchange cannot be repaid in the terms it revised, and it does so without making uniqueness a condition of autonomy.

Responsiveness is necessary and not sufficient for good supervision. An agent can be responsive and wrong about what should be revised, and a revision that improves one agent's contribution can disrupt the coordination on which others rely. A clerk who reports a defective form is doing something different from a clerk who decides that the form should not exist, and the criterion does not distinguish them; it says only that both are answerable to what bears on the matter. Whether a given revision improves the contribution is a further question, answered by its consequences, which is the second axis of the grid taken up next.

The passage, like the others, belongs to a posthumous compilation and is written in a polemical register. That register is separable from its structural point. The moral contempt for the ethos of exchange is one thing, and the observation that a system of equivalence presupposes interchangeable agents is another; the essay adopts only the latter, and the analysis does not depend on the former. The point does not condemn reciprocity among those who stand in a position to give and receive on equal terms. It identifies a cost of treating persons as fungible by default, which is the loss of precisely that part of an action that no ledger records.

\section{Authorship and Function}

The analysis so far concerns the relation of an agent to its purposes. It does not concern what the activity does in the world, and the two questions must be kept apart. An independently chosen activity may produce nothing but deterioration, as vandalism does, while an externally assigned one may genuinely preserve a valuable structure. A sanitation worker does not have to invent the purpose of sanitation for the work to maintain the conditions under which a city remains habitable. The work maintains the constraints that support a valued organization, here the habitability of a city, and this is the sense of contragrade change in Deacon's framework \cite{deacon2012}. It is not simply the reversal of entropy production, since the accumulation of refuse need not raise local thermodynamic entropy in any useful sense; the relevant distinction concerns the maintenance of constraints and not the direction of an entropy balance. Whether a given activity has this character is a fact about its consequences for the structure it acts on, and it is independent of whether the performer authored, endorsed, or can revise its purpose.

The two questions therefore define a grid and not a scale. An activity can be self-owned and constructive, as with the builder of a house to their own design. It can be self-owned and destructive, as with the vandal. It can be assigned and constructive, as with the sanitation worker who has taken up the task or been given it without reflection. And it can be assigned and idle or destructive, as with a bureaucratic routine that maintains nothing but itself. The Nietzschean concern of this essay lies on the first axis, and it should not be read as a verdict on the second. A person with a fully examined set of purposes can do harm, and a person who has examined nothing can maintain a great deal. What the first axis adds is a further question about the second: whether the person doing constructive work could recognize, and respond to, the point at which it ceased to be constructive. This essay confines itself to the first axis and notes that a complete account of work would need both.

\section{Four Cases}

Four ordinary cases show how the criterion applies and how little the surface of the activity decides. The first is the apprentice in a skilled trade, who spends years laying tile, hanging doors, or fitting pipe. The repetition is real and, at first, tedious, and it gradually becomes automatic, freeing attention for the properties of the material that a novice cannot notice: the way a substrate absorbs moisture, the way a wall departs from plumb, the early sign that a method will fail in cold weather. Here encapsulation and responsiveness grow together, and the trade is a clear instance of automaticity serving autonomy. The same apprentice becomes insulated when the shop dictates one method for every condition and treats a question about why as insubordination. The retrospective test is what happened when a product, a code, or a site condition changed; the design test is whether the shop explains its methods; the structural test is what the shop would do if its standard method were shown to fail.

The second case is the learner of a language who works through a deck of vocabulary cards. The deck is a requital ledger in miniature, converting cleared cards into standing, and its fixed order can quietly become an insulating device, because a learner who always meets the cards in sequence may recall a word from its neighbor and not from its meaning. Shuffling the deck removes the cue and forces recall to depend on the word itself, which is a small instance of keeping an interface open to the relevant input and closing it to an irrelevant one. The learner who notices that a memorized gloss does not match how the word is used, and revises the card, is exercising the supervisory capacity at the smallest scale. The learner who finishes decks without ever meeting the language in use has acquired a reliable procedure with no check on whether it serves communication.

The third case is the player who grinds in an online game. The loop, meter, and unlock belong to the designer, and the player's diligence is what the design converts into engagement. The case is useful for the question of what happens when the ground fails. If the designer retires the item, the player who stops, or who redirects effort to a project the player values, has treated the grind as a means. The player who immediately seeks the next loop has shown that the loop itself was the point. A third kind of player, who builds tools, modifies the game, or sets challenges of their own devising, has moved the ownership of the meter, which is the transformation the earlier discussion of the self-owned grind describes.

The fourth case is the clerk or official, the figure Nietzsche's own list names. Routine administrative work is a mixed case, since a well-run office needs reliable procedure and cannot survive constant relitigation of its rules. The question is whether the officer remains answerable to the purposes the procedure serves. An officer who notices that a form has stopped capturing what it was designed to capture, and reports it, is supervising the procedure. An officer who processes the form because it is the form is a fungible occupant of a position. The two may be indistinguishable on any ordinary day, and the difference appears only when the form fails, which is why the criterion is retrospective and structural and not diagnostic of individuals in advance.

The cases also show that responsiveness is relative to the level at which a purpose is specified. A learner may become expert at correcting cards without asking whether those words are worth learning, and an officer may report a defective form without questioning the objective the form serves. Supervision of a procedure is not supervision of the project that the procedure serves, nor is supervision of a project supervision of the institution that makes the project desirable. This is the structure of protagonist and executor described earlier, repeated at each level: the same agent can be protagonist with respect to a procedure and executor with respect to a project. Local procedural intelligence is therefore not evidence of the broader independence that the essay associates with personality, and a claim of personality has to specify the level at which the purposes are held answerable. The criterion applies at every level, and independence in the strong sense is a matter of how far up the levels the capacity extends and how readily it can be exercised there.

\section{Objections and Replies}

Four objections deserve a direct reply. The first is that all expertise rests on repetition, so that an argument which treats drilled routines with suspicion condemns the means of excellence. The reply is that the argument does not condemn repetition or automaticity, which it treats as enabling. Its target is a particular failure, in which higher-order judgment is subordinated to the routine, and the trades and arts are the standard cases in which repetition and judgment develop together.

The second objection is that Nietzsche denies the unified subject, so that talk of an agent supervising its purposes presupposes the very self he dissolves. In the compilation the subject is called a fiction, and the ego is called an imagined synthesis that is not one with the unitary controlling force of our being \cite[Aphs.\ 370--71]{wtp1}. The same book describes a weak will as the multiplicity and disintegration of the instincts and a strong will as their coordination under the government of one among them \cite[Aph.\ 46]{wtp1}. Because these passages come from the compilation, the objection should not be given more weight than a posthumous note bears, and it is stronger when pressed from the published works, where the same doubt about a simple unified subject recurs \cite[Secs.\ 12, 17, and 19]{bge}. The reply is that supervision as described here does not require a simple inner unit. The standing disposition to keep a routine answerable to what bears on it is a feature of how an organism's drives and habits are organized, and it can come in degrees without a homunculus to exercise it. This is also the reason the essay declines to posit a supervisory module. What Nietzsche denies is the unity of the subject taken as a given; what he praises under the heading of personality is a unity achieved by organization, and the essay's criterion concerns that achievement.

The third objection is that any society needs fungible role-fillers, so that an argument against fungibility is utopian. The reply is that the argument does not ask that roles be abandoned or that every worker reinvent the purpose of the work. A sanitation worker, a clerk, and a nurse fill roles whose functions are specified in advance, and the world is better for their reliability. What the argument asks is that the occupants retain the capacity to notice when the role has stopped serving its function, and that institutions be designed so that this capacity is cultivated and heard. Interchangeability of occupants is compatible with answerability of occupants, and the combination is what a well-run institution looks like.

The fourth objection is that responsiveness can be luck or display, so that the criterion cannot separate the supervisor from the lucky or the theatrical. The reply is the one already given in the operational definition: display is distinguished by its insensitivity to what would count against it, and luck by its failure to recur across varied circumstances. Neither distinction can be applied at a glance, which is why the criterion yields comparative and retrospective judgments and not a classification of individuals.

\section{Limits and Qualifications}

Two limits should be stated. The first is that responsiveness to reasons is a criterion and not a method of detection. From the outside, and often from the inside, it is difficult to know how one would respond to a change in circumstances that has not occurred. Self-ascribed responsiveness is cheap, and the person who is certain that they would walk away, or certain that they would hold firm, is frequently mistaken. What the criterion licenses is three things: a retrospective judgment about how a person or an institution in fact responded when a premise failed, a design question about whether an arrangement keeps its grounds open to inspection and rewards reconsideration, and a structural question about what an arrangement would do were its premise removed. It does not license classifying a particular individual in advance.

The second limit concerns the elitism of the source. Nietzsche's remark about personality belongs to a view in which a few individuals justify the rest, and the argument developed here does not depend on that view. The structural claim, that a certain relation to one's purposes is an achievement distinct from competence, can be stated without the order-of-rank doctrine, and in one respect the argument requires its rejection. If only a few persons were capable of responsiveness to reasons about their purposes, the institutional question of whether schools, workplaces, and games can cultivate it would be moot, and the criterion would sort people and not arrangements. The practical upshot of the essay depends on responsiveness being cultivable in the many, which is a stronger commitment than declining to endorse a hierarchy, and one the essay presupposes and does not establish. Discipline of the kind Foucault described as producing docile and useful bodies \cite{foucault1977} is a general feature of modern institutions and not the artifact of any one author's hierarchy, and whether such discipline leaves room for reconsideration of its own ends can be asked without presupposing that only a few are fit to ask it.
\section{Conclusion}

The passage on the philologist as educator and the contemporary vocabulary of the grind describe a common arrangement, in which the repetition of unrewarding acts forms a disposition to continue irrespective of interest. Nietzsche's separate remark on mechanical activity shows that such an arrangement can serve several functions, the formation of the young, the relief of the exhausted, and the practice of the skilled, which cannot be told apart by observing the activity. His observation that a personal act is rare, read as a claim about the agent's relation to its purposes and not about behavioral richness, suggests what separates them, and the figure of the non-player character gives the separation a portable form: an agent may be rich in behavior and fixed in purpose, and may be a protagonist at one level and an executor at the next. Personality is not shown by choosing one's routines, nor by dropping them when their rewards vanish, but by retaining the capacity to examine and reorganize the purposes that make those routines intelligible, and by responding to reasons in doing so. The criterion applies at several levels, those of procedure, project, and institution, and competence at a lower level is not independence at a higher one. It also bears on exchange: a system of equivalents presupposes interchangeable occupants, whereas an agent who supervises a routine can revise the terms of performance, a contribution that no ledger of equivalents can register. That capacity does not require originating one's purposes, only appropriating them. The philologist's magnificent monotony is accordingly neither simply a vice nor simply a virtue. It is a powerful instrument whose significance depends on whether those who undergo it retain, or are helped to acquire, the capacity to ask what it is for, and whether what it accomplishes in the world is worth the asking.

\newpage
\begin{thebibliography}{99}

\bibitem{deacon2012}
Deacon, Terrence W. \emph{Incomplete Nature: How Mind Emerged from Matter}. New York: W. W. Norton, 2012.

\bibitem{frankfurt1971}
Frankfurt, Harry G. ``Freedom of the Will and the Concept of a Person.'' \emph{Journal of Philosophy} 68, no.\ 1 (1971): 5--20.

\bibitem{foucault1977}
Foucault, Michel. \emph{Discipline and Punish: The Birth of the Prison}. Translated by Alan Sheridan. New York: Pantheon Books, 1977.

\bibitem{hughes1989}
Hughes, John. ``Why Functional Programming Matters.'' \emph{Computer Journal} 32, no.\ 2 (1989): 98--107.

\bibitem{kant1784}
Kant, Immanuel. ``An Answer to the Question: What Is Enlightenment?'' 1784. In \emph{Practical Philosophy}, edited and translated by Mary J. Gregor, 11--22. Cambridge: Cambridge University Press, 1996.

\bibitem{montinari2003}
Montinari, Mazzino. \emph{Reading Nietzsche}. Translated by Greg Whitlock. Urbana: University of Illinois Press, 2003.

\bibitem{parnas1972}
Parnas, D. L. ``On the Criteria to Be Used in Decomposing Systems into Modules.'' \emph{Communications of the ACM} 15, no.\ 12 (1972): 1053--58.

\bibitem{bge}
Nietzsche, Friedrich. \emph{Beyond Good and Evil: Prelude to a Philosophy of the Future}. Translated by Walter Kaufmann. New York: Vintage Books, 1966.

\bibitem{wtp1}
Nietzsche, Friedrich. \emph{The Will to Power: An Attempted Transvaluation of All Values}. Vol.\ 1, Books I and II. Translated by Anthony M. Ludovici. Edited by Oscar Levy. Edinburgh: T. N. Foulis. Project Gutenberg eBook 52914.

\bibitem{wtp2}
Nietzsche, Friedrich. \emph{The Will to Power: An Attempted Transvaluation of All Values}. Vol.\ 2, Books III and IV. Translated by Anthony M. Ludovici. Edited by Oscar Levy. Edinburgh: T. N. Foulis, 1913. Project Gutenberg eBook 52915.

\bibitem{wp}
Nietzsche, Friedrich. \emph{We Philologists}. Translated by J. M. Kennedy. Edited by Oscar Levy. Edinburgh: T. N. Foulis, 1911. Project Gutenberg eBook 18267.

\bibitem{suits1978}
Suits, Bernard. \emph{The Grasshopper: Games, Life and Utopia}. Toronto: University of Toronto Press, 1978.

\bibitem{yee2006}
Yee, Nick. ``The Labor of Fun: How Video Games Blur the Boundaries of Work and Play.'' \emph{Games and Culture} 1, no.\ 1 (2006): 68--71.

\end{thebibliography}

\end{document}
